\subsection{Key IRL Methodologies}

Inverse Reinforcement Learning (IRL) addresses a fundamental problem: given observed agent behavior, how can we recover the reward function that motivated it? Formally, given an MDP $\mathcal{M} \setminus R = \{S, A, P, \gamma\}$ without the reward function, and demonstrations $\mathcal{D} = \{\tau_1, \tau_2, ..., \tau_m\}$ where each trajectory $\tau_i = \{(s_0, a_0), (s_1, a_1), ..., (s_T, a_T)\}$ consists of state-action pairs, the goal is to recover the reward function $R: S \times A \rightarrow \mathbb{R}$ that the demonstrator is optimizing \citep{NgRussell2000}.

\subsubsection{Max-Margin IRL}

\citet{NgRussell2000} introduced the max-margin approach, establishing that a reward function should make the expert's policy better than alternatives by some margin. For all states $s \in S$ and actions $a \neq \pi_E(s)$, it enforces constraints:

$$(P_{\pi_E(s)} - P_a)(I - \gamma P_{\pi_E})^{-1}R \geq 0$$

\noindent where $P_a$ is the transition matrix for action $a$, $P_{\pi_E}$ is the transition matrix under the expert policy, $\gamma$ is the discount factor, and $R$ is the reward vector. The objective maximizes this margin while keeping rewards bounded:

$$\max_R \sum_{s\in S} \min_{a\neq\pi_E(s)} [(P_{\pi_E(s)} - P_a)(I - \gamma P_{\pi_E})^{-1}R] - \lambda||R||_1$$

\noindent where $\lambda$ controls the reward sparsity. \citet{Ratliff2006} extended this approach with Maximum Margin Planning, formulating a quadratic program that makes the expert's behavior optimal under the learned reward. This approach requires known transition dynamics and assumes expert optimality, limiting its applicability in environments with uncertain dynamics \citep{Abbeel2004}.

\subsubsection{Maximum Entropy IRL}

\citet{Ziebart2008} introduced a probabilistic approach using the principle of maximum entropy to resolve ambiguity in IRL. Rather than assuming deterministic optimality, MaxEnt IRL models the expert's trajectory distribution as:

$$P(\tau|\theta) = \frac{1}{Z(\theta)} \exp(\theta^T \sum_{t=0}^T \phi(s_t))$$

\noindent where $\tau$ is a trajectory, $\theta$ are reward parameters, $\phi(s_t)$ are state features, and $Z(\theta)$ is the partition function. This formulation is equivalent to assuming the agent follows a Boltzmann policy, mathematically corresponding to a dynamic discrete choice model with logit errors \citep{Ermon2015}. The objective becomes maximum likelihood estimation of reward parameters, solved via gradient ascent:

$$\nabla L(\theta) = \tilde{\phi} - \sum_s D_s \phi_s$$

\noindent where $\tilde{\phi}$ represents empirical feature counts from demonstrations and $D_s$ are expected state visitation frequencies. \citet{Ziebart2008} demonstrated this approach's effectiveness for predicting pedestrian movement, achieving 74.8\% accuracy compared to 45.9\% for shortest-path predictions. \citet{Wulfmeier2015} later extended this to deep neural network reward representations, enabling applications in environments with high-dimensional features.

\subsubsection{Generative Adversarial Imitation Learning}

\citet{Ho2016} reframed IRL as a distribution matching problem using adversarial training. GAIL employs a discriminator network $D_\omega: S \times A \rightarrow [0,1]$ to distinguish between expert and learner behaviors, while a policy network $\pi_\theta: S \rightarrow P(A)$ tries to fool the discriminator:

$$\min_\theta \max_\omega \mathbb{E}_{\tau \sim \pi_\theta}[\log(1 - D_\omega(s,a))] + \mathbb{E}_{\tau \sim \pi_E}[\log(D_\omega(s,a))] - \lambda H(\pi_\theta)$$

\noindent where $H(\pi_\theta)$ is the entropy regularization term. The discriminator effectively learns a reward function: $\log(D_\omega(s,a)) - \log(1 - D_\omega(s,a))$. \citet{Fu2018} extended this to Adversarial IRL, which explicitly recovers a reward function that generalizes beyond training trajectories, demonstrating improved transfer to modified environments in robotic tasks. \citet{Kostrikov2018} further refined this approach with Discriminator Actor-Critic, improving sample efficiency through off-policy training.

{\bf Summary of Identifiability Conditions in IRL.} Inverse Reinforcement Learning (IRL) is generally underdetermined, because multiple reward functions can induce the same observed policy \citep{kim2021irl}. For instance, any constant shift on the reward is unidentifiable. Nonetheless, recent theoretical work has established conditions for unique (up to a constant) identifiability. First, if we assume deterministic MDPs combined with maximum-entropy regularization, observed trajectory frequencies directly reveal the reward \citep{kim2021irl}. In more general settings, coverage assumptions ensure every state-action lies on some optimal trajectory, which provides enough constraints for uniqueness \citep{kim2021irl}. Another key line of work uses multiple environments (or discount factors): consistent optimal behavior in diverse MDPs forces a single reward function across them all \citep{cao2021irl,rolland2022_identifiability}. If the reward is linear in known features and the constant function is excluded from the feature span, then the parameter vector becomes uniquely identifiable \citep{rolland2022,shehab2024}. Finally, simplifying the reward structure---for example, restricting to state-only or time-homogeneous rewards---removes degrees of freedom and thus promotes uniqueness \citep{cao2021irl}. Overall, these approaches illustrate how additional assumptions or data sources can make IRL well-posed, with a single consistent reward function (up to additive constants). 

\subsection{Identification Challenges and Econometric Connections}

Different IRL methods address identification through various mechanisms. Max-Margin IRL selects the reward function with maximum margin between optimal and suboptimal actions \citep{Ratliff2006}. MaxEnt IRL applies maximum entropy as a selection principle, providing a unique solution consistent with observed behavior frequencies \citep{Ziebart2008}. \citet{kim2021irl} formalized necessary and sufficient conditions for reward identifiability, proving that specific properties of the environment and policy determine whether rewards can be uniquely recovered.

\citet{OBriant2024} demonstrates that standard IRL’s use of “wind shocks” yields only partial identification, because random deviations cannot reveal relative payoffs. Substituting preference shocks from the DDC framework of \citet{Rust1987} restores full identification. Under this assumption, \citet{Abbeel2004}’s projection IRL attains nearly the accuracy of NFXP but requires up to 20× fewer dynamic programming solves. By contrast, the linear program of \citet{NgRussell2000} remains computationally simpler yet restricts payoffs to \(\{-1, 0, 1\}\).

The connection between IRL and structural econometrics is mathematically precise. MaxEnt IRL's probabilistic formulation is equivalent to a dynamic discrete choice model with logit errors \citep{Ermon2015}. The Boltzmann policy corresponds to assuming i.i.d. Gumbel noise in utilities, leading to choice probabilities matching multinomial logit specifications. This equivalence was validated by \citet{Ermon2015}, who showed that MaxEnt IRL solutions coincide with maximum-likelihood estimates from corresponding structural models.

\citet{Sharma2018} operationalized this connection by adapting Hotz-Miller's conditional choice probability (CCP) approach from econometrics to IRL. Their CCP-IRL approach achieved a 5× speed improvement with no loss in reward quality compared to standard MaxEnt IRL on benchmark tasks. This computational advance addresses a key challenge in both fields: the nested fixed-point computation required in each iteration. In IRL, each reward update typically requires solving a forward MDP—analogous to the nested fixed-point algorithms in structural estimation \citep{Rust1987}.

\citet{Zeng2022} proposed a single-loop algorithm that interleaves policy optimization with reward parameter updates, proving convergence to stationary points with theoretical guarantees. This approach parallels simultaneous estimation methods in econometrics, both avoiding full policy convergence at each iteration. \citet{Boularias2011} developed Relative Entropy IRL, a model-free approach using KL-divergence that avoids requiring known transition dynamics, making it applicable to settings where the environment model is unavailable.

\subsection{Economic Applications of IRL}

IRL methods have been applied to domains that resemble economic decision-making, demonstrating potential to complement structural econometric modeling in environments.

\subsubsection{Dynamic Discrete Choice and Migration Decisions}

\citet{Ermon2015} applied IRL to model pastoralist farmers' movement decisions in East Africa using GPS-tracked movements. Their study treated migration routes as optimal policies in an MDP where rewards reflected preferences for resource abundance versus travel costs. The model quantified trade-offs between water availability, grazing quality, and distance traveled. The authors found that drought conditions would significantly alter movement patterns, with herders traveling 21\% farther on average to reach water sources—a finding with implications for climate adaptation policy.

This application highlights how IRL can handle high-dimensional state spaces and dynamics that would challenge traditional structural estimation. \citet{Bogota2015} applied similar techniques to refugee migration patterns, demonstrating how IRL can recover preference parameters from observational data even when underlying utility functions are non-linear.

\subsubsection{Route Choice and Transportation Economics}

IRL has been used to infer driver preferences in route selection, complementing discrete choice models in transportation economics. \citet{Ziebart2008} applied MaxEnt IRL to predict taxi navigation paths, revealing that drivers optimize reward functions beyond simple distance minimization. \citet{Bronner2023} utilized IRL to detect anomalous ride-hailing driver routes, identifying preferences for specific road types and time-dependent routing strategies with 89\% accuracy.

\citet{Mai2015} compared IRL approaches to traditional route choice models, finding that IRL captured path dependencies and strategic anticipation that multinomial logit models missed. Their analysis showed that IRL-derived models reduced prediction error by 17\% when forecasting traffic redistribution after infrastructure changes, highlighting IRL's value for transportation policy evaluation.

\subsubsection{Labor Market and Skill Development}

\citet{Yancey2022} applied IRL to model worker skill development and career path choices, treating labor market participation as a sequential decision problem. By analyzing career trajectory data from 15,000 workers over a ten-year period, the study estimated that workers valued skill development opportunities at approximately 14-18\% of immediate compensation, varying by education level. The resulting model predicted responses to unemployment benefit changes with higher accuracy than traditional structural labor models.

\citet{Chan2019} extended this approach to incorporate bounded rationality in job search, showing how IRL can recover discount factors and risk preferences simultaneously with reward functions. Their findings indicated significant heterogeneity in how workers value job security versus wage growth potential, with implications for designing effective unemployment insurance programs.

These applications demonstrate IRL's ability to recover preference parameters from observed behavior in dynamic settings, enabling counterfactual analysis and policy evaluation—core objectives of structural economic analysis. The field continues to develop, with recent work by \citet{Ramachandran2007} and \citet{Finn2016} introducing Bayesian IRL and Guided Cost Learning approaches that further enhance IRL's applicability to economic modeling.
